% 鸟神星（综述）
% license CCBYSA3
% type Wiki

本文根据 CC-BY-SA 协议转载翻译自维基百科\href{https://en.wikipedia.org/wiki/Makemake}{相关文章}。

\begin{figure}[ht]
\centering
\includegraphics[width=6cm]{./figures/abca2436367ca8e9.png}
\caption{2015 年 4 月，哈勃空间望远镜拍摄的低分辨率阋神星及其无名卫星 S/2015 (136472) 1 图像} \label{fig_Makema_1}
\end{figure}
玛克马克$^{[g]}$（小行星编号：136472 Makemake）是一颗位于柯伊伯带的矮行星，柯伊伯带是海王星轨道之外由冰质天体组成的圆盘。它是第四大海王星外天体，也是经典柯伊伯带中体积最大的成员$^{[h]}$，其直径约为冥王星的 60\%。它于 2005 年 3 月 31 日由美国天文学家迈克尔·E·“迈克”·布朗（Michael E. "Mike" Brown）、查德·特鲁希略（Chad Trujillo）和大卫·拉比诺维茨（David Rabinowitz）在帕洛马天文台发现。作为该团队发现的最大天体之一，玛克马克的发现促成了 2006 年冥王星被重新归类为矮行星。

玛克马克在表面特征方面与冥王星类似：其表面高度反射，大部分覆盖着冻结的甲烷，并因凝聚物（tholins）而被染成红褐色$^{[i]}$。玛克马克目前已知有一颗卫星，但尚未命名。该卫星的轨道表明玛克马克自转具有很高的轴倾角，这意味着它会经历极端季节。玛克马克显示出地球化学活动与低温火山活动的证据，这使科学家怀疑其可能拥有一个地下液态水海洋。在玛克马克上发现了气态甲烷，但尚不清楚它是存在于大气中，还是来自暂时性的气体逸出。

由于尚未有太空探测器近距离访问玛克马克，因此并不存在关于其表面的高分辨率影像。玛克马克距离地球非常遥远，即使通过望远镜观测，它看起来也只是一颗类似星点的光斑。科学家表达了希望派遣太空探测器访问玛克马克的愿望，因为它具有地质活动及潜在的地下海洋。
\subsection{历史}
\subsubsection{发现}
\begin{figure}[ht]
\centering
\includegraphics[width=8cm]{./figures/7d2c811983023908.png}
\caption{玛克马克是在帕洛马天文台使用 1.22 米塞缪尔·奥辛望远镜（如图）拍摄的图像中被发现的} \label{fig_Makema_2}
\end{figure}